\documentclass[11pt]{article}
\usepackage[T1]{fontenc}
\usepackage{microtype}
\usepackage[margin=1in]{geometry}
\usepackage{amsmath,amssymb,amsthm}
\usepackage[hidelinks]{hyperref}

\newtheorem{theorem}{Theorem}
\newtheorem{proposition}[theorem]{Proposition}
\newtheorem{corollary}[theorem]{Corollary}
\DeclareMathOperator{\SO}{SO}
\newcommand{\Gm}{\mathbb G_m}
\newcommand{\CT}{\operatorname{CT}}
\newcommand{\Lq}{\Lambda_q}

\title{The Mathieu property for quadratic complements}
\author{Roy van Rijn}
\date{}
\hypersetup{pdftitle={The Mathieu property for quadratic complements},pdfauthor={Roy van Rijn}}

\begin{document}
\maketitle

\begin{abstract}
Let \(q\) be a nondegenerate quadratic form of dimension \(N\ge1\) over a
field of characteristic zero. The kernel of the normalized
\(\SO(q)\times\Gm\)-invariant trace on \(\{q\ne0\}\) is a Mathieu--Zhao
subspace exactly when \(N\le2\). For \(N\ge3\) we give an explicit
geometrically nonpower function whose positive powers have zero trace,
while multiplication by \(q^{N-2}\) gives a nonzero trace at every order.
For this function the least detecting polynomial degree is \(2N-4\).
We give the sharp vanishing cutoff below that degree and a finite test
for eventual vanishing with any fixed regular multiplier.
\end{abstract}

\section{A double-pole counterexample in dimension three}

Take \(q=x^2+y^2+z^2\) over \(\mathbb Q\) and let \(\Lq\) be the normalized
trace on \(\mathbb Q[x,y,z,q^{-1}]\) invariant under rotations and scalar
multiplication. The function
\[
   F_0=\frac{(q+1)\bigl(q-(q+1)^2z^2\bigr)}{q^2}
\]
satisfies, for every \(m\ge1\),
\begin{equation}\label{eq:opening}
   \Lq(F_0^m)=0,\qquad
   \Lq(qF_0^m)=\int_0^1(1-w^2)^m\,dw\ne0 .
\end{equation}
The first marked moment is \(2/3\), and the same multiplier \(q\)
works at every order. The identities hold in any characteristic-zero
field, since the integral is a positive rational number.

For a nondegenerate form \(q\) on a \(K\)-vector space \(V\), write
\[
 X_q=\{q\ne0\},\qquad \mathcal O(X_q)=K[V,q^{-1}],
 \qquad G_q=\SO(q)\times\Gm .
\]
The two factors act by isometries and scalar multiplication.
The trace \(\Lq\) is the unique \(G_q\)-invariant linear functional
taking \(1\) to \(1\). Call \(F\) \emph{pure} if \(\Lq(F^m)=0\) for
every \(m\ge1\). The kernel is a Mathieu--Zhao subspace if, for every
pure \(F\) and every fixed multiplier \(h\in\mathcal O(X_q)\), the
marked moments \(\Lq(hF^m)\) vanish for all sufficiently large \(m\).

\begin{theorem}\label{thm:main}
Let \(K\) have characteristic zero and \(q\) be nondegenerate of
dimension \(N\ge1\). Then
\[
   \boxed{\ker\Lq\text{ is a Mathieu--Zhao subspace}
          \quad\Longleftrightarrow\quad N\le2.}
\]
For \(N\ge3\) there is an explicit pure \(F\in\mathcal O(X_q)\) for
which \(\Lq(q^{N-2}F^m)\ne0\) for every \(m\ge1\). The function \(F\)
is not a unit times a proper power over an algebraic closure.
\end{theorem}

In dimensions one and two the split trace is a Laurent constant term.
For \(N\ge3\), take an orthogonal splitting \(q=A+B\) of ranks \(2\)
and \(M=N-2\). The ratio \(B/q\) has moments \(M/(M+2j)\).
Scalar invariance then reduces the mixed trace to a Laurent
coefficient of a polynomial integral. At the endpoint of that
integral, the pure moment selects a coefficient above the degree
of a polynomial, and multiplication by \(q^M\) selects the leading
coefficient. The same calculation identifies the least polynomial
degree needed to detect a failure for this \(F\).

Over \(\mathbb C\), Long's spherical classification already gives the
dimension cutoff and the existence of fixed detecting multipliers
\cite[Thm.~1.1]{Long}. To see that it applies, write
\(q=2xy+q_0(z)\) for \(N\ge3\), with \(b_0\) the symmetric bilinear
form of \(q_0\). The isometries
\[
 (x,y,z)\longmapsto
 \bigl(x,\,y-b_0(z,u)-xq_0(u)/2,\,z+xu\bigr)
 \qquad(u\in\mathbb C^{N-2})
\]
form a unipotent group normalized by the torus scaling \(x,y\)
inversely. Together with the scalar torus they form a connected
solvable subgroup of \(G_q\). On \(x\ne0\) they can set \(z=0\),
\(q=1\), and \(x=1\), so a Borel subgroup has an open orbit in \(X_q\).
The stabilizer of a norm-one vector is reductive. The scalar factor in
the stabilizer is \(\pm1\), and its identity component is the special
orthogonal group of the vector's orthogonal complement. Since the
derived subgroup contains the rotations, it acts nontrivially on
\(X_q\) for \(N\ge3\). For \(N\le2\), \(G_q\) is a torus.
The same complex cutoff follows from the general criterion in
\cite[Thm.~1.1]{vanRijn}.

The endpoint substitution here adapts Long's calculation on
\(\mathrm{SL}_2/N(T)\) \cite[Prop.~5.1]{Long}. The rotation recurrence
supplies the moments over the original characteristic-zero field.
It gives explicit witnesses there, as well as the exact cutoff for
their polynomial multipliers.

\section{The trace and the two low dimensions}\label{sec:positive}

The action on \(X_q\) is transitive over an algebraic closure.
Scalar multiplication changes a nonzero value of \(q\) to any other
such value. In dimension at least two, an isometry taking one vector
to another of the same nonzero norm exists by extending an isometry
of their nondegenerate lines. Its determinant can be corrected by
a reflection in the orthogonal complement of the target vector.
In dimension one the scalar action alone is transitive.

The coordinate ring is a locally finite rational representation
of the connected reductive group \(G_q\). Complete reducibility
\cite[\S12l and Thm.~22.42]{Milne} gives the projection onto its trivial
isotypic part. This part is \(K\cdot1\) by geometric transitivity,
which proves existence and uniqueness of \(\Lq\). The trace commutes
with field extension. Indeed, the extended trace is still invariant
and normalized, and the same uniqueness applies over the larger
field. Invariance also implies \(\Lq(Df)=0\) for every infinitesimal
action \(D\).

We use the Laurent constant-term theorem of Duistermaat and van der
Kallen \cite[\S\S1, 3]{DvK}. If a nonzero complex Laurent polynomial
has zero constant term in every positive power, then the convex hull
of its support excludes the origin. This implication holds over any
characteristic-zero field. The finitely many coefficients generate
a finitely generated extension of \(\mathbb Q\), which embeds in
\(\mathbb C\) by choosing images for a transcendence basis and
extending algebraically. The embedding preserves the support and
all the constant-term identities. A linear functional strictly
positive on that support gives a lower bound \(m\delta\), for some
\(\delta>0\), on the weights of monomials in the \(m\)-th power.
If the least monomial weight of a fixed nonzero multiplier is \(b\),
the product has weights at least \(b+m\delta\). Hence its constant
term eventually vanishes. If either factor is zero the conclusion is
immediate.

When \(N=1\), write \(q=ax^2\). The ring is \(K[x,x^{-1}]\), and the
scalar action gives \(\Lq=\CT_x\) since each \(x^i\) has weight \(i\).
When \(N=2\), pass to a splitting field \(L/K\) of degree at most two
and write \(q=a uv\) in linear coordinates \(u,v\). Both \(u\) and \(v\)
are then units on the complement. The rotation and scalar tori act
on \(u^iv^j\) with character \(t^{i-j}s^{i+j}\), which is trivial
exactly for \(i=j=0\). Hence \(\Lq\otimes_K L=\CT_{u,v}\) on
\(L[u^{\pm1},v^{\pm1}]\). For \(f,h\) defined over \(K\), pure vanishing
extends to \(L\), and the eventual mixed vanishing over \(L\) descends
to \(K\) by injectivity. This proves the positive direction of
Theorem~\ref{thm:main}.

\section{A rank-two splitting and its moments}\label{sec:beta}

Assume \(N\ge3\) and set \(M=N-2\). Diagonalize \(q\) over \(K\) and write
\[
   q=A+B,\qquad A=a_1x_1^2+a_2x_2^2,\qquad
   B=\sum_{i=1}^M b_i z_i^2,\qquad S=\frac{B}{q}.
\]
All \(a_\alpha,b_i\) are nonzero. To compute the moments of \(S\), use the
infinitesimal rotations
\[
   D_{\alpha i}
   =b_i z_i\frac{\partial}{\partial x_\alpha}
    -a_\alpha x_\alpha\frac{\partial}{\partial z_i},
   \qquad \alpha=1,2,\quad 1\le i\le M.
\]
They preserve \(q\), and hence act on its localization. For \(j\ge1\),
the product rule gives
\begin{equation}\label{eq:rotation}
   \sum_{\alpha,i}D_{\alpha i}
   \left(S^{j-1}\frac{x_\alpha z_i}{q}\right)
     =2S^j-(M+2j-2)(1-S)S^{j-1}.
\end{equation}
Differentiating \(x_\alpha z_i\) contributes
\(2S^j-M(1-S)S^{j-1}\). For \(j>1\), differentiating \(S^{j-1}\)
contributes \(-2(j-1)AB S^{j-2}/q^2
=-2(j-1)(1-S)S^{j-1}\). That term is absent when \(j=1\).

Applying \(\Lq\) to \eqref{eq:rotation} yields
\[
   (M+2j)\Lq(S^j)=(M+2j-2)\Lq(S^{j-1}),
   \qquad \Lq(1)=1.
\]
It follows that \(\Lq(S^j)=M/(M+2j)\) for all \(j\ge0\), or
equivalently
\begin{equation}\label{eq:beta}
   \Lq(g(S))=M\int_0^1 t^{M-1}g(t^2)\,dt\qquad(g\in K[s]).
\end{equation}
For coefficients in \(K\), the integral denotes termwise polynomial
integration, defined by \(\int_0^1t^r\,dt=1/(r+1)\) and \(K\)-linearity.
Also, \(q^\ell g(S)\) has scalar weight \(2\ell\), so its trace is
zero whenever \(\ell\ne0\). This allows us to compute the trace of
any Laurent expression in \(q\) with polynomial coefficients in \(S\)
by applying \eqref{eq:beta} to its constant coefficient.

\section{The endpoint calculation and the witness}\label{sec:endpoint}

Put \(W=1-(1+q)^2 S\). For integers \(\alpha,\gamma\ge1\) and
\(\beta\ge0\), define
\[
 F_{\alpha,\beta,\gamma}=q^{-\alpha}(1+q)^\beta W^\gamma,\qquad
 C_{k,j}=M\int_0^1 w^{M+2j-1}(1-w^2)^k\,dw
 \quad(k,j\ge0).
\]
The constants \(C_{k,j}\) are positive rationals.

\begin{proposition}\label{prop:moments}
For \(m\ge1\), \(j\ge0\), and \(\ell\ge(\alpha-\gamma)m\),
\begin{equation}\label{eq:moments}
 \Lq(q^\ell S^jF_{\alpha,\beta,\gamma}^{\,m})
    =C_{\gamma m,j}\binom{\beta m-M-2j}{\alpha m-\ell}.
\end{equation}
Here \(\binom uv=[U^v](1+U)^u\), with value zero for \(v<0\). At
\(\ell=(\alpha-\gamma)m-1\), the left-hand side minus the binomial
expression is \(M(-2)^{\gamma m}/(\gamma m+1)\).
\end{proposition}

\begin{proof}
Let \(U\) be the formal variable for \(q\), and set \(T=1+U\),
\(k=\gamma m\), and \(n=\alpha m-\ell\).
Scalar weights and \eqref{eq:beta} give
\[
 \begin{aligned}
 \Lq(q^\ell S^jF_{\alpha,\beta,\gamma}^{\,m})
 &= [U^n]\,T^{\beta m}M\int_0^1 t^{M+2j-1}(1-T^2t^2)^k\,dt\\
 &= [U^n]\,T^{\beta m-M-2j}J(T),\\
 J(T)&=M\int_0^T w^{M+2j-1}(1-w^2)^k\,dw .
 \end{aligned}
\]
The substitution \(w=Tt\) is valid term by term.
The derivative \(J'(T)\) has initial term \(M(-2)^k(T-1)^k\) at
\(T=1\). Thus \(J(1+U)-J(1)\) starts with
\(M(-2)^k U^{k+1}/(k+1)\). The factor \(T^{\beta m-M-2j}\)
has constant coefficient \(1\) at \(U=0\), so the error first
appears at \(n=k+1\). For \(n\le k\) the retained constant
\(J(1)=C_{k,j}\) gives \eqref{eq:moments}.
\end{proof}

Take
\begin{equation}\label{eq:uniform}
 F=F_{M,M,M+1}=\left(\frac{1+q}{q}\right)^M W^{M+1}.
\end{equation}
With \(j=0\), the upper index in \eqref{eq:moments} is \(Mm-M\).
The choices \(\ell=0,M\) give, respectively,
\begin{equation}\label{eq:witness}
 \Lq(F^m)=0,\qquad \Lq(q^M F^m)=C_{(M+1)m,0}.
\end{equation}
The same calculation with \(M=1\) and \(F_0=F_{1,1,1}\) proves
\eqref{eq:opening}, and applies to every nondegenerate ternary form
using a one-dimensional summand \(B\).

To check the power assertion, work over an algebraic closure.
Along the punctured \(x_1\)-axis, \(W=1\), and \(1+q\) has a simple
zero at a point with \(q=-1\). The restriction of
\(F_{\alpha,\beta,\gamma}\) there has vanishing order \(\beta\).
At a point with \(q=1\) and \(B=1/4\), \(W=0\) and \(dW=-4\,dB\).
One can fix \(A=3/4\) and vary just \(z_1\) through a nonzero root of
\(b_1z_1^2=1/4\), so \(W\) has a simple zero along that line.
The restriction of \(F_{\alpha,\beta,\gamma}\) has order \(\gamma\)
there. A unit restricts to a unit on both lines where \(q\ne0\).
Any proper power up to a unit therefore has exponent dividing
both \(\beta\) and \(\gamma\). For \(F\) these are \(M,M+1\),
and for \(F_0\) they are \(1,1\). This completes
Theorem~\ref{thm:main}.

The divisor \(q=0\) is geometrically irreducible since a reducible
quadratic has rank at most two. The numerator of \(W\) is \(-B\)
modulo \(q\), so \(W\) has a simple pole there. The pole order of
\(F_{\alpha,\beta,\gamma}\) is \(\alpha+\gamma\). This gives
\(2M+1\) for \(F\) and two for the ternary \(F_0\).

\begin{corollary}\label{cor:scalar-extension}
For quadratic complements the Mathieu--Zhao property of the trace
kernel is invariant under every field extension. When it fails,
the same phase and multiplier witness the failure over the
extension.
\end{corollary}

\begin{proof}
Dimension and nondegeneracy are unchanged by extension. The trace
commutes with scalar extension and the marked moments in
\eqref{eq:witness} are nonzero rational numbers.
\end{proof}

\section{The least degree of a detecting multiplier}\label{sec:cutoff}

\begin{proposition}\label{prop:average}
For a phase \(G\) invariant under the block isometries
\(O(A)\times O(B)\), every regular multiplier \(h\) has a block
average \(\pi h\in K[q,q^{-1},S]\) such that
\begin{equation}\label{eq:average}
 \Lq(hG^m)=\Lq((\pi h)G^m)\qquad(m\ge1).
\end{equation}
If \(h\) is polynomial of degree at most \(D\), the average is a
linear combination of \(q^\ell S^j=q^{\ell-j}B^j\) with
\(0\le j\le\ell\le\lfloor D/2\rfloor\).
\end{proposition}

\begin{proof}
Block isometries preserve \(\Lq\) by uniqueness, since they normalize
the rotations and commute with scalars. For a polynomial \(h\), take
its Reynolds projection in the space of polynomials of degree at
most \(\deg h\). Complete reducibility for the identity component,
followed by averaging over the finite component group, gives this
projection \cite[Cor.~22.43]{Milne}. The linear functional
\(h\mapsto\Lq(hG^m)\) is block invariant, which proves \eqref{eq:average}.
For \(h=q^{-r}p\) with \(p\) polynomial, use the same projection on
\(q^{-r}\) times the space of polynomials of degree at most \(\deg p\).

The polynomial invariants of the block group are \(K[A,B]\).
Over an algebraic closure choose unit vectors \(v_A,v_B\) in the two
blocks. Every pair on \(AB\ne0\) is equivalent under block isometries
to \((s v_A,t v_B)\), with \(s^2=A,\ t^2=B\).
An invariant polynomial restricts to an even polynomial in each of
\(s,t\), by the sign symmetries. Thus it is a polynomial in \(A,B\)
on a dense locus. Algebraic independence of \(A,B\) gives descent of
its coefficients to \(K\). The degree filtration is preserved, and
rewriting \(A=q-B\) gives the stated monomials.
\end{proof}

\begin{theorem}\label{thm:cutoff}
For the phase \(F\) in \eqref{eq:uniform} and the double-pole ternary
phase \(F_0\), the least degree of a polynomial multiplier detecting
infinitely many orders is \(2M\). For every integer \(0\le D<2M\),
the least positive \(m_0\) for which all polynomial multipliers of
degree at most \(D\) vanish at every order \(m\ge m_0\) is
\[
 m_0=1+\left\lceil\frac{2\lfloor D/2\rfloor}{M}\right\rceil .
\]
\end{theorem}

\begin{proof}
Set \(J=\lfloor D/2\rfloor\). By \eqref{eq:average} it suffices to
consider \(q^\ell S^j\) with \(0\le j\le\ell\le J\).
In \eqref{eq:moments}, the upper index \(Mm-M-2j\) is
nonnegative for \(m\ge m_0\). The lower index \(Mm-\ell\) exceeds
it by \(M+2j-\ell>0\). This proves the cutoff. If \(J>0\), take
\(h=B^J=q^J S^J\) at \(m=\lceil2J/M\rceil\). The upper index is
strictly negative and the lower one is positive, so the moment is
nonzero. For \(J=0\), the claimed cutoff is already \(1\).
Finally \(q^M\) has degree \(2M\) and detects all orders by
\eqref{eq:witness}, or by \eqref{eq:opening} for \(F_0\).
\end{proof}

For example, for the uniform phase in dimension five, a multiplier
needs degree at least six to detect infinitely many orders. All
polynomial multipliers of degree at most five vanish from order
three, and \(B^2\) detects the second order.

\paragraph{A linear multiplier in the isotropic ternary case.}
Write \(q=2xy+az^2\), with \(a\in K^\times\), and set
\[
 p=(1+x)^3\frac{y}{q}-\frac{(1+x)(2+x)}2 .
\]
Then for every \(m\ge1\),
\begin{equation}\label{eq:isotropic-linear}
 \Lq(p^m)=0,\qquad
 \Lq(xp^m)=2^{-m}\int_0^1(1-w^2)^m\,dw\ne0 .
\end{equation}
Indeed the hyperbolic torus gives opposite weights to \(x\) and
\(u=y/q\). The derivation \(az\,\partial_y-x\,\partial_z\) applied
to \(x^{k-1}y^kz/q^k\) gives
\((2k+1)\Lq((xu)^k)=k\Lq((xu)^{k-1})\). Thus on polynomials in
\(x,u\) the trace equals the constant term in \(x\) after substituting
\(u=(1-t^2)/(2x)\), integrated over \(0\le t\le1\).
This follows by checking the moments with equal exponents, for
which \(2^{-k}\int_0^1(1-t^2)^k\,dt\) satisfies the same recurrence.
With \(T=1+x\) the substitution gives
\(p=T(1-T^2t^2)/(2x)\). For the two moments in
\eqref{eq:isotropic-linear}, the coefficients sought have degrees
\(m\) and \(m-1\) in
\(2^{-m}T^{m-1}\int_0^T(1-w^2)^m\,dw\).
The integral differs from its value at \(T=1\) first in order
\(m+1\) in \(x\). The retained polynomial \(T^{m-1}\) has zero
coefficient at degree \(m\) and leading coefficient \(1\).

On the line \(y=1,z=0\), the restriction \(p=(1+x)/(2x)\) has a
simple zero at \(x=-1\), so \(p\) is geometrically nonpower up to
units. A constant multiplier gives zero by purity. Hence the
least detecting polynomial degree among all witnesses on an
isotropic ternary complement is one. For the double-pole phase
\(F_0\) on the same complement, Theorem~\ref{thm:cutoff} gives two.

\section{A test for a fixed regular multiplier}\label{sec:tail}

Even above the degree threshold, terms of a multiplier can cancel.
The following test accounts for those cancellations.

\begin{proposition}\label{prop:tail}
Let \(F\) be as in \eqref{eq:uniform}, put \(\gamma=M+1\), and
write \(\pi h=\sum_{\ell,j}c_{\ell,j}q^\ell S^j\) for a fixed
regular multiplier \(h\). Put \(a=M/2\). For all sufficiently large \(m\),
\[
  \Lq(hF^m)=C_{\gamma m,0}R_h(m),\qquad
  R_h(X)=
    \sum_{\ell\ge M+2j} c_{\ell,j}
    \frac{(a)_j}{(a+\gamma X+1)_j}
    \binom{MX-M-2j}{\ell-M-2j}.
\]
Here \((u)_j\) is the rising factorial, with \((u)_0=1\), and the
binomial in \(R_h\) is a polynomial in \(X\).
The marked moments eventually vanish exactly when \(R_h=0\)
as a rational function. Otherwise they are nonzero at every
sufficiently large order.
\end{proposition}

\begin{proof}
The finite support of \(\pi h\) lies in the range of
\eqref{eq:moments} once \(m\) is large. The upper indices are then
nonnegative. Terms with \(\ell<M+2j\) vanish, and for the other terms
one can use the complementary binomial index \(\ell-M-2j\).
Integration of the derivative of \(w^{M+2j}(1-w^2)^{k+1}\) gives
\[
 (M+2j)C_{k,j}=(M+2j+2k+2)C_{k,j+1}.
\]
Consequently \(C_{k,j}/C_{k,0}=(a)_j/(a+k+1)_j\).
This proves the formula. The constants \(C_{\gamma m,0}\) are
nonzero and a nonzero rational function has only finitely many zeros.
\end{proof}

The block average needed in the test can be found by linear algebra.
On \(q^{-r}\) times the polynomials of degree at most \(D\), where
\(r\) clears the denominator of \(h\), take the quotient by the images
of the block rotation derivations and of one coordinate reflection
minus the identity in each block. The rotation images remove the
nontrivial simple summands for the identity component. The two
reflection differences remove the remaining nontrivial components
of the finite group. The unique invariant representative in the
quotient is \(\pi h\). Clearing denominators in \(R_h\) then gives
a polynomial identity test over \(K\).

For the uniform ternary phase \(F=F_{1,1,2}\), the multiplier
\(4q^4S+15q^3S-q\) has a nonzero first moment, yet its marked moments
vanish for all \(m\ge3\). Its rational test is
\((4(X-3)+15)/(4X+3)-1=0\), with the cancellation at \(m=3\)
given directly by \eqref{eq:moments}.

\section{Pole order and another witness}\label{sec:poles}

\begin{proposition}\label{prop:product}
Within the family \(F_{\alpha,\beta,\gamma}\) with \(\gamma\ge\alpha\),
purity holds exactly when \(\beta\ge M\) and \(\alpha\ge\beta\).
For a pure phase in this family, there exists a regular multiplier
detecting infinitely many orders exactly when \(\alpha=\beta\)
or \(\alpha=\gamma\). The least pole order among the geometrically
nonpower witnesses in this family is two for \(M=1\) and
\(2M+1\) for \(M\ge2\).
\end{proposition}

\begin{proof}
Set \(\ell=j=0\) in \eqref{eq:moments}. If \(\beta<M\), the
upper index at \(m=1\) is negative and the coefficient is nonzero.
For \(\beta\ge M\), purity requires and follows from
\(\alpha\ge\beta\). If \(\alpha=\beta\), \(q^M\) detects every order.
If \(\gamma>\alpha>\beta\), every term of an averaged regular
multiplier eventually falls within the range of \eqref{eq:moments}
and has lower index above the nonnegative upper one.
For \(\gamma=\alpha>\beta\), \(q^{-1}\) picks the first coefficient
past that range, so its moment is exactly
\(M(-2)^{\alpha m}/(\alpha m+1)\) at every order. In this case all
polynomial multipliers eventually vanish, since their averaged terms
have \(\ell\ge0\).

For a pure phase, \(\gamma\ge\alpha\ge\beta\ge M\), so its pole
order is at least \(2M\). Equality forces all three exponents to
be \(M\). This gives the ternary phase for \(M=1\), and a proper
power for \(M\ge2\). The phase in \eqref{eq:uniform} attains the next
integer \(2M+1\).
\end{proof}

A polynomial perturbation gives pole order \(2M\) in dimensions at
least six, with the same minimum multiplier degree.

\begin{proposition}\label{prop:perturb}
For \(M\ge4\), the phase
\[
   F_*=F_{M,M,M}\bigl(1+B(M-1+Mq)\bigr)
\]
is pure, geometrically nonpower up to a unit, and has pole order
\(2M\). The multiplier \(q^M\) detects every order, and \(2M\) is the
least polynomial degree detecting infinitely many orders.
\end{proposition}

\begin{proof}
In the \(m\)-th power of the added factor the monomials are
\(B^j q^r=q^{j+r}S^j\), with \(0\le r\le j\le m\), and their
coefficients are \(\binom mj\binom jr(M-1)^{j-r}M^r\).
For \(m\ge2\), the upper index \(Mm-M-2j\) in \eqref{eq:moments}
is nonnegative, since it is at least \((M-2)m-M\).
The pure lower index exceeds it by \(M+j-r\).
For multiplier \(q^M\) the difference is \(j-r\), so precisely
the terms with \(j=r\) contribute, all positively. At \(m=1\)
the two added pure terms cancel because
\[
  (M-1)\binom{-2}{M-1}+M\binom{-2}{M-2}=0.
\]
Both added terms have negative lower index for the \(q^M\)-marked
moment at \(m=1\), leaving \(C_{M,0}\).

Modulo \(q\), the added factor has constant term \(1\), so the pole
order stays \(2M\). Over an algebraic closure it has a simple zero
on the smooth norm level \(q=1\) at a point with \(B=-1/(2M-1)\).
Indeed \(A\) and \(B\) are nonzero there, so their differentials are
independent and \(dB\) is nonzero on the tangent space to \(q=1\).
The added factor has differential \((2M-1)dB\) on that tangent
space, and the base phase is nonzero at the point. The resulting
simple zero excludes a proper power up to a unit.

For a polynomial multiplier of degree \(D<2M\), its averaged terms
are \(q^\ell S^{j_0}\) with \(j_0\le\ell\le\lfloor D/2\rfloor\).
Together with \(B^j q^r\) in the power expansion this makes the
lower index exceed the upper one by \(M+2j_0-\ell+j-r>0\).
The upper index is nonnegative once
\((M-2)m\ge M+2\lfloor D/2\rfloor\). Thus every such multiplier
eventually vanishes, while \(q^M\) detects all orders.
\end{proof}

\paragraph{AI assistance and review.}
ChatGPT (OpenAI) assisted with mathematical arguments, literature checks,
exposition, proof review and verification code.

\begin{thebibliography}{9}
\bibitem{DvK}
J.~J. Duistermaat and W. van der Kallen,
\emph{Constant terms in powers of a Laurent polynomial},
Indag. Math. (N.S.) \textbf{9} (1998), no.~2, 221--231.
\href{https://doi.org/10.1016/S0019-3577(98)80020-7}{DOI: 10.1016/S0019-3577(98)80020-7}.

\bibitem{Long}
C.~D. Long,
\emph{A Classification of Mathieu's Second Conjecture for Reductive
Spherical Pairs}, 2026.\par\nopagebreak
GitHub \href{https://github.com/long-mathematics/mathieu-second-conjecture-spherical-pairs}{\nolinkurl{long-mathematics/mathieu-second-conjecture-spherical-pairs}},
\href{https://github.com/long-mathematics/mathieu-second-conjecture-spherical-pairs/blob/e9edac262b4b64eb68feb0e6e80496cefe4a1dee/mathieu_second_conjecture_spherical_pairs.tex}{rev.~\texttt{e9edac2}}.

\bibitem{Milne}
J.~S. Milne,
\emph{Algebraic Groups: The Theory of Group Schemes of Finite Type over a
Field}, Cambridge Studies in Advanced Mathematics \textbf{170},
Cambridge University Press, 2017.

\bibitem{vanRijn}
R. van Rijn,
\emph{Counterexamples to the Mathieu Property: From Explicit Quotient
Constructions to a General Classification}, 2026.
\href{https://doi.org/10.5281/zenodo.24310836}{DOI: 10.5281/zenodo.24310836}.
\end{thebibliography}

\end{document}
